% Grundlagen — claude-tunnel bachelor thesis.
\chapter{Grundlagen}
\label{ch:grundlagen}

Dieses Kapitel fuehrt die Konzepte ein, auf denen der Entwurf aufbaut: das
Prinzip der Providerblindheit (Abschnitt~\ref{sec:providerblind}), den
Transport \gls{ac:quic} (\ref{sec:quic}), die Ende-zu-Ende-Kryptographie mit dem
Noise-Protokoll (\ref{sec:noise}), Proof-of-Work als Missbrauchsschranke
(\ref{sec:pow}) und die Traversierung von \gls{ac:nat} (\ref{sec:nat}).

\section{Providerblindheit und das Zero-Knowledge-Prinzip}
\label{sec:providerblind}

Ein Vermittlungsdienst heisst in dieser Arbeit \gls{gl:providerblind}, wenn er
den vermittelten Verkehr strukturell nicht einsehen kann. Der Begriff
\emph{Zero-Knowledge} wird hier im woertlichen, systemtechnischen Sinne
gebraucht -- der Betreiber besitzt \enquote{null Wissen} ueber Ziel und Nutzlast
-- und nicht im kryptographischen Sinne eines Zero-Knowledge-Beweises. Die
Blindheit ist eine \emph{strukturelle} Eigenschaft: sie ergibt sich daraus, dass
die vertraulichen Daten den Vermittler nur als Chiffretext erreichen und die
Adressierung ueber opake Tokens statt Klarnamen erfolgt. Damit unterscheidet sie
sich von einer bloss \emph{organisatorischen} Zusicherung (etwa einer
No-Log-Richtlinie), die auf dem Wohlverhalten des Betreibers beruht.

Das Vertrauensmodell verschiebt sich entsprechend: Der Client muss dem Betreiber
weder Vertraulichkeit noch Nichtprotokollierung glauben, weil dieser die
relevanten Informationen gar nicht erst erhaelt. Er muss ihm lediglich
\emph{Verfuegbarkeit} zutrauen -- ein Relay kann den Dienst verweigern oder
verzoegern, aber nicht mitlesen. Anonymisierungsnetze wie
Tor~\cite{dingledine2004tor} verfolgen ein verwandtes Ziel ueber mehrere
Zwischenstationen; der hier gewaehlte Ansatz konzentriert sich stattdessen auf
die Blindheit eines einzelnen Vermittlers plus einen direkten Pfad, der diesen
bei Erreichbarkeit ganz umgeht.

\section{QUIC als Transport}
\label{sec:quic}

\gls{ac:quic}~\cite{rfc9000} ist ein verbindungsorientiertes Transportprotokoll
oberhalb von \gls{ac:udp}. Es integriert die Verschluesselung fest in den
Transport: der Handshake nutzt \gls{ac:tls}~1.3~\cite{rfc8446,rfc9001}, sodass
bereits die Transportschicht authentisiert und verschluesselt ist. Gegenueber
\gls{ac:tcp}+TLS bietet \gls{ac:quic} drei fuer diese Arbeit relevante
Eigenschaften: (i) \emph{gemultiplexte Streams} innerhalb einer Verbindung ohne
Head-of-Line-Blocking zwischen ihnen, (ii) einen \emph{schnellen
Verbindungsaufbau} (1-RTT, mit Sitzungswiederaufnahme 0-RTT) und (iii)
Verbindungs-IDs, die einen Pfadwechsel (Connection Migration) erlauben.

Die Stream-Multiplexierung traegt im Entwurf den Rollen-Dispatch: Client, Agent
und Steuerbefehle teilen sich eine Verbindung, ohne sich gegenseitig zu
blockieren. Weil \gls{ac:quic} auf \gls{ac:udp} aufsetzt, ist es zugleich der
Punkt, an dem Zensur ansetzen kann -- wird \gls{ac:udp} blockiert, ist
\gls{ac:quic} nicht nutzbar. Der Entwurf braucht daher einen Transport-Fallback
(Abschnitt~\ref{sec:nat} und Kapitel~\ref{ch:architektur}), der denselben Tunnel
ueber \gls{ac:tls}-ueber-\gls{ac:tcp} traegt.

\paragraph{Handshake und Datenphase.} Der Verbindungsaufbau laeuft in
\gls{ac:quic} ueber \emph{Initial}- und \emph{Handshake}-Pakete, in die der
\gls{ac:tls}-1.3-Handshake eingebettet ist; nach einer Rundreise (1-RTT) sind die
Schluessel etabliert und die Datenphase (\emph{1-RTT}-Pakete) beginnt. Bei einer
zuvor bekannten Gegenstelle erlaubt die Sitzungswiederaufnahme fruehe Daten
(0-RTT), die jedoch gegen Wiedereinspielung (Replay) abzusichern sind und daher
nur fuer idempotente Anfragen taugen. Streams sind entweder bidirektional oder
unidirektional und werden unabhaengig flusskontrolliert; Verlust auf einem Stream
blockiert die anderen nicht (kein stream-uebergreifendes Head-of-Line-Blocking).
Erst dies macht den Rollen-Dispatch (Kapitel~\ref{ch:architektur}) auf einer
einzigen Verbindung effizient: Registrierung, Relay und Steuerbefehle konkurrieren
nicht um dieselbe Sequenz.

\section{Ende-zu-Ende-Kryptographie mit Noise}
\label{sec:noise}

Damit der Vermittler blind bleibt, genuegt die Transportverschluesselung von
\gls{ac:quic} nicht: sie endet am Edge. Die Nutzlast muss zusaetzlich zwischen
Client und Origin verschluesselt sein. Dafuer dient das Noise Protocol
Framework~\cite{perrin2018noise}, ein Baukasten fuer authentisierte
Schluesselaustausch-Handshakes aus Diffie-Hellman-Operationen, einem
AEAD-Chiffre und einer Hashfunktion.

Konkret wird das Muster \texttt{Noise\_IK\_25519\_ChaChaPoly\_BLAKE2s}
eingesetzt: Schluesselaustausch ueber Curve25519 (X25519), Verschluesselung mit
ChaCha20-Poly1305, Hashing mit BLAKE2s. Im Muster \texttt{IK} uebertraegt der
\emph{Initiator} (Client) seinen statischen Schluessel waehrend des Handshakes
(\enquote{I}), waehrend der statische Schluessel des \emph{Responders} (Origin)
dem Initiator vorab bekannt ist (\enquote{K}). Genau diese Vorabkenntnis ist der
Hebel fuer die Providerblindheit: der Client \emph{pinnt} die Origin-Identitaet
(den oeffentlichen Origin-Schluessel), die er ausserhalb des Bandes als Teil
einer Capability erhaelt. Ein Vermittler, der die Origin-Identitaet nicht kennt,
kann den Handshake nicht als Man-in-the-Middle unterlaufen; ein Angreifer mit
falschem Schluessel scheitert am Pinning. Nach dem Handshake gehen beide Seiten
in den Transport-Modus ueber und tauschen laengen-praefigierte, verschluesselte
Frames aus, die der Edge nur noch byteweise weiterreicht.

\paragraph{Ablauf des IK-Handshakes.} Der Handshake besteht aus zwei Nachrichten.
In der ersten sendet der Initiator seinen ephemeren oeffentlichen Schluessel und
-- bereits verschluesselt -- seinen statischen Schluessel; die
Diffie-Hellman-Operationen \texttt{es} (ephemeral--static) und \texttt{ss}
(static--static) mischen dabei sowohl den ephemeren als auch den vorab bekannten
Origin-Schluessel in den Sitzungszustand, sodass nur der echte Origin die
Nachricht entschluesseln kann. In der zweiten Nachricht antwortet der Responder
mit seinem ephemeren Schluessel und den Operationen \texttt{ee} und \texttt{se};
damit ist der Sitzungsschluessel beidseitig authentisiert und -- dank der
ephemeren Anteile -- \emph{vorwaertsgeheim}: die spaetere Kompromittierung eines
statischen Schluessels legt vergangene Sitzungen nicht offen. Anschliessend leiten
beide Seiten je Richtung einen symmetrischen Transportschluessel ab. Weil der
Origin-Schluessel bereits in die erste Nachricht eingeht, ist der Handshake an die
gepinnte Identitaet gebunden, noch bevor Nutzdaten fliessen.

\section{Proof-of-Work als Missbrauchsschranke}
\label{sec:pow}

Ein blinder Dienst kann Missbrauch nicht am Inhalt erkennen. Um die
Verbindungsanbahnung dennoch zu verteuern, wird ein Proof-of-Work nach dem
Hashcash-Prinzip~\cite{back2002hashcash} vorgeschaltet. Der Edge stellt eine
Challenge aus einem Nonce und einer Schwierigkeit; der Client muss eine Loesung
finden, deren Hash eine geforderte Anzahl fuehrender Nullbits aufweist. Das
Finden ist -- exponentiell in der Schwierigkeit -- teuer, die Pruefung dagegen
ein einzelner Hash. Dieser asymmetrische Aufwand begrenzt Flut- und
Rendezvous-Missbrauch, ohne dass der Dienst Identitaeten fuehren oder Inhalte
sehen muesste. Die Schwierigkeit ist ein Stellparameter, dessen Einfluss die
Parameterstudie (Kapitel~\ref{ch:evaluation}) untersucht.

\paragraph{Kosten und Asymmetrie.} Modelliert man den Hash als Zufallsorakel, so
trifft ein einzelner Versuch die geforderten $d$ fuehrenden Nullbits mit
Wahrscheinlichkeit $2^{-d}$; die erwartete Zahl der Versuche bis zur Loesung ist
damit $2^{d}$, waehrend die Pruefung genau einen Hash kostet. Der Aufwand des
Clients waechst also exponentiell in der Schwierigkeit $d$, der des Edge bleibt
konstant. Genau diese Asymmetrie macht den Nachweis als Schranke brauchbar: Eine
einzelne Anbahnung bleibt fuer einen ehrlichen Client guenstig, waehrend das
massenhafte Anbahnen eines Fluters quadratisch bis exponentiell teurer wird. Die
Wahl von $d$ ist ein Kompromiss zwischen wirksamer Daempfung und der Belastung
schwacher Clients (Kapitel~\ref{ch:diskussion}).

\section{NAT-Traversal und direkter Pfad}
\label{sec:nat}

Reale Endpunkte liegen haeufig hinter \gls{ac:nat}, das ausgehende Verbindungen
zulaesst, eingehende aber ohne Zustand verwirft. Der Relay-Pfad umgeht dies, weil
sowohl Client als auch Agent \emph{ausgehend} zum Edge verbinden. Ein
\emph{direkter} \gls{ac:p2p}-Pfad, der den Edge umgeht, muss dagegen die
\gls{ac:nat}-Bindung ueberwinden. Das etablierte Verfahren ist das
gleichzeitige Oeffnen (Hole Punching) auf Basis reflexiver, vom Vermittler
beobachteter Adressen~\cite{ford2005peer}; standardisiert ist die
Kandidaten-Aushandlung als \gls{ac:ice}~\cite{rfc8445}. Im Entwurf inseriert der
Agent einen direkten Listener beim Edge, den der Client entdeckt und bevorzugt
nutzt; scheitert der direkte Aufbau, faellt der Client transparent auf den Relay
zurueck. Echtes Hole Punching ueber symmetrisches \gls{ac:nat} bleibt eine
Grenze des flachen Container-Testbetts und wird in
Kapitel~\ref{ch:evaluation} eingeordnet.

\subsection{NAT-Typen und ihre Bedeutung fuers Hole Punching}
\label{sec:nat-typen}

\gls{ac:nat}-Geraete unterscheiden sich darin, wie sie eine einmal geoeffnete
ausgehende Bindung fuer eingehenden Verkehr wiederverwenden. Bei einem
\emph{Full-Cone}-\gls{ac:nat} gilt die Bindung fuer beliebige externe Absender;
bei \emph{Address-Restricted}- bzw.\ \emph{Port-Restricted}-\gls{ac:nat} nur fuer
zuvor kontaktierte Adressen bzw.\ Adress-Port-Paare. In diesen drei Faellen
genuegt das gleichzeitige Oeffnen: beide Seiten senden zuerst hinaus und erzeugen
so die noetige Bindung fuer die Gegenrichtung, woraufhin die vom Vermittler
beobachtete reflexive Adresse den Peer erreicht. Ein \emph{symmetrisches}
\gls{ac:nat} dagegen vergibt je Ziel eine \emph{andere} externe Portnummer; die
beobachtete reflexive Adresse sagt dann nichts ueber den Port voraus, den der
Peer tatsaechlich sehen wird, und das Hole Punching schlaegt ohne zusaetzliche
Portvorhersage fehl. Der vorliegende Prototyp setzt auf reflexive Adressen und
deckt damit die nicht-symmetrischen Faelle ab; symmetrische \gls{ac:nat}s bleiben
eine bekannte Grenze (Kapitel~\ref{ch:diskussion}).
